% Discussion & limitations — full prose v1.

\paragraph{What one cell in 178 means.}
We do not conclude that prefix-locked value is impossible --- our single
prefix-limited cell is an existence proof --- but that it is rare enough
in this regime that any account of reasoning built on frequent
within-trace unlocks is measuring budget fit. The practical corollary of
the compression finding is a reframing of restart policies: for a
compute-starved model, abandoning a long prefix typically costs a budget
multiple rather than access to the solution; for a capability-limited model,
neither continuation nor restart helps. Which regime a deployment faces is
measurable with a four-point restart curve.

\paragraph{What the prediction results do and do not show.}
Three layers --- a pre-registered null on our cohort, a trace-blind
difficulty ceiling inside the published range on 192K public R1
generations, and a
replicate-then-dissect analysis at 256 samples per problem --- support one
conclusion: pooled probe AUROCs on reasoning outcomes are difficulty
measurements until shown otherwise. This does not overturn any specific
published number. \citet{zhang2025knowright} probe intermediate-answer
positions and their data contain within-trace label variance; our claim
there is a falsifiable prediction: adding a question-only baseline and a
within-problem evaluation will substantially shrink the reported gaps. It
also does not make internal signals useless: within-problem trace
\emph{selection} \citep{fu2025deepconf} and pre-generation difficulty
routing \citep{lugoloobi2026} are consistent with everything we find ---
both exploit exactly the difficulty signal we isolate. Together with the
causal failures of probe signals \citep{yuan2026diagnostic} and their
coherence confound under contaminated prefixes \citep{pair2026}, the
critique now stands on three independent legs: probe signals are causally
inert, coherence-driven, and difficulty-driven.

\paragraph{Limitations.}
Our instrumented cohort uses two small models in 4-bit quantization on one
benchmark family, with $m = 4$--$8$ attempts per cell and one base seed;
the confirmatory baseline is text-level (a pre-generation activation probe
would be stronger \citep{lugoloobi2026}); the public-dump dissection
re-scores full-precision generations through a 4-bit model (top-1 fidelity
0.906, reported) and inherits unknown sampling temperature; the probe
dissection itself is mathematics-only (the difficulty ceiling generalizes
off-math --- GPQA-diamond, Section~\ref{sec:results-prediction} --- but a
non-math dissection awaits a multi-sample dump from a locally runnable
model); and the
within-trajectory \emph{timing} question --- when does value arrive? ---
remained under-powered everywhere because genuine interior events are
rare, which is itself a finding: at this scale, breakthroughs are either
immediate or absent. No equivalence margin was pre-registered for the
confirmatory endpoints,
so their results are absences of detected gain rather than demonstrated
equivalence; and positive existence claims (the one prefix-limited cell,
the intermediate-solvability band) rest on modest counts and warrant
replication.

\paragraph{Future work.}
The natural next experiment runs the same frozen instrument on a
large RL-trained reasoner of the kind the aha-moment claims describe: interior
events there would revive the timing question with the controls already in
place, and their absence would extend the budget-artifact account to the
regime that motivated it. A second direction treats the value curve
$\phat(t)$ as the object of study --- its shape (gradual versus stepped),
not its crossings --- for which the present data already suffice.
